\section{Introduction}\label{sec:intro}

Solar-flare prediction from photospheric magnetogram time series is a
benchmark-driven machine-learning problem. The community's public
reference is SWAN-SF~\cite{angryk2020}, and the literature on it has
grown around a small number of evaluation conventions: randomised
splits of pooled windows, a cleaned and rebalanced export, and a
reporting culture built on ROC/PR summary statistics. Those
conventions carry risk. The benchmark's own creators warned that
sliding-window temporal coherence ``invalidates the randomness
assumption, thus impacting all sampling
practices''~\cite{angryk2019}; subsequent work has echoed the
warning~\cite{ahmadzadeh2021}, and the cleaned export's release paper
explicitly prescribes pairing temporally-preceding training
partitions with a later test partition~\cite{eskandari2024}.
Flare-prediction models built on photospheric vector-magnetogram
summaries --- the SHARP parameter suite~\cite{bobra2015} --- have
demonstrated skilful classification of flaring versus non-flaring
active regions~\cite{leka2019,georgoulis2021}; the predictive signal
lives in quantities such as total unsigned current helicity,
magnetic free energy, and flux-emergence proxies.

Why federate a public benchmark whose data is freely poolable?
Methodologically, cross-silo federated
learning~\cite{mcmahan2017,kairouz2021} is the standard instrument
for studying non-IID data heterogeneity under a controlled protocol:
it prescribes exactly what moves between data holders (model
parameters) and exactly what does not (raw observations), which makes
it a clean lens for asking how much of a benchmark's apparent
difficulty survives when the data is deliberately fragmented.
Operationally, the SHARP-style parameters at the centre of this
pipeline are dominated by a single public instrument (SDO/HMI via
JSOC); the six-regional-agency narrative sometimes attached to
federated space-weather proposals does not describe this data and is
not claimed here. The six clients of this study are a simulation
protocol --- Dirichlet label skew across disjoint shards --- not a
claim about any institution. Federated flare forecasting itself is
not new (a federated transfer-learning preprint exists~\cite{fu2023});
what has not existed is a federated treatment of SWAN-SF that audits
its own benchmark artefacts before interpreting its numbers.

This paper provides that treatment, and its character is an audit.
We train six simulated clients on the benchmark in both its
cleaned/rebalanced and raw, naturally-imbalanced forms, under a
frozen evaluation protocol whose every selection decision is made on
validation data, and we ask three questions: what does a federated
pipeline actually achieve on SWAN-SF, what do the benchmark's own
artefacts do to the answer, and which parts of the observed behaviour
belong to the algorithms rather than to the evaluation machinery?
The contributions are:

\begin{enumerate}[leftmargin=2em, itemsep=3pt]
  \item \textbf{An instance-level provenance audit of the cleaned
  SWAN-SF export.} Normalisation-invariant matching against the raw
  benchmark (98.7--100\% coverage, 100\% test-side label agreement)
  shows the cleaned export's same-partition train/test pairing
  shares instances: every flaring test window (6{,}234 of 6{,}234) is
  in the training pool, and 85.7--90\% of training positives are
  synthetic (Section~\ref{sec:data}). The leakage phenomenon is
  documented prior art~\cite{angryk2019,eskandari2024}; the
  instance-level quantification inside a federated study is new.
  \item \textbf{A leakage-free federated evaluation.} On the
  temporally-preceding fold (train partitions 1--4, single-pass test
  on partition 5), all arms generalise (ROC-AUC 0.906--0.978) and the
  in-partition federated-versus-centralised gap disappears; event-level
  detection is essentially solved for every arm, and false-alarm
  burden --- not detection --- separates them
  (Section~\ref{sec:results}).
  \item \textbf{Corrected persistence baselines as the deployability
  verdict.} Same-region label inertia reaches TSS 0.967 with no
  learned model; only XGBoost (0.848) and logistic regression
  (0.489) clear 24-hour-lagged persistence (0.418) on TSS, and no arm
  clears it on HSS (Section~\ref{sec:results}).
  \item \textbf{Two BatchNorm findings that re-attribute apparent
  federated failures.} The in-partition FedAvg collapse (0.575) is a
  composite evaluation artefact --- untransported normalisation
  statistics compounded by a saturated validation monitor and the
  checkpoint rule it feeds; and on the raw, unbalanced substrate, an
  architecture-only control (three BatchNorm layers replaced by the
  identity, nothing else changed) restores the federated MLP to the
  centralised model's ROC-AUC level, so the collapse is dominated by
  BatchNorm under federation, not by the flattened-feature encoder
  and not by federation per se (Section~\ref{sec:results}).
  \item \textbf{A calibration failure that is real, not an
  artefact.} Pooled-statistics recalibration --- validated as an
  instrument on the centralised model, where it reproduces the
  model's own statistics to within $7\times10^{-5}$ ROC-AUC ---
  collapses every federated arm to chance on the raw substrate, so a
  pooled-statistics transport fix would not rescue federation there;
  and no validation-fit decision layer survives the $26\times$
  prevalence shift between the balanced training partitions and the
  natural-prevalence test (Section~\ref{sec:calibration}).
  \item \textbf{A frozen, reproducible protocol.} Immutable splits,
  validation-only selection and calibration, single-pass test
  evaluation, and machine-readable artefacts for every number
  reported here, released with the code.
\end{enumerate}

We make no operational alerting claim. The benchmark's data is
public and freely poolable; the downstream
flare$\to$CME$\to$storm$\to$GIC chain is physically far from the
quantity modelled and is not evaluated (see~\cite{boteler2019} for
that literature).
